\section{SCENARIO \#15: "Cedar Creek'': October 19, 1864
}


\begin{scenariodata}
\textbf{Game Length:} 1 Turn\\
\end{scenariodata}

\begin{scenariodata}
(Weather: July - October Column)\\
\end{scenariodata}

\begin{scenariodata}
\textbf{First Phase:} Confederates\\
\end{scenariodata}

\begin{scenariodata}
\textbf{Union:}\\
\end{scenariodata}

\begin{scenariodata}
Resource Factor 217\\
\end{scenariodata}

\begin{scenariodata}
\textbf{Initial Placement:}\\
\end{scenariodata}

\begin{scenariodata}
Hex N-25 : 191, 3A, (Crook)\\
\end{scenariodata}

\begin{scenariodata}
Hex N-24 : 241, 2A, (Emory, 1F, 1S)\\
Hex M-23 1171, 4A, (Wright, 1S, 1W)\\
Hex N-23 : O1, (2S, 2W)\\
\end{scenariodata}

\begin{scenariodata}
Hex M-22 : 161, 3A, (1S, 1W)\\
\end{scenariodata}

\begin{scenariodata}
Hex M-21 : 31, (1S, 1W)\\
\end{scenariodata}

\begin{scenariodata}
Hex N-20 (Winchester) : 31, (Sheridan, 1F, 5S, 3W)\\
Hex A-16 (Romney) -- : [41, 1A], (1F, 1S)\\
\end{scenariodata}

\begin{scenariodata}
Hex L-21 :8C, 3H\\
Hex K-20 : 7C, 2H, (1S, 1W)\\
\end{scenariodata}

\begin{scenariodata}
Hex Q-25 :3C, 1H\\
\end{scenariodata}

\begin{scenariodata}
B\&O : (201, 5A, 3C], (4F, 4S, 1W)\\
\end{scenariodata}

\begin{scenariodata}
\textbf{Withdrawal Hexes:}\\
\end{scenariodata}

\begin{scenariodata}
N-7, 8-19, Z-13, CC-15 through CC-33\\
\end{scenariodata}

\begin{scenariodata}
(Supply/Wagon Entry Hexes:\\
\end{scenariodata}

\begin{scenariodata}
N-7, \$-10, Z-13, A-11, A-13, A-16, CC-15 through CC-33\\
\end{scenariodata}

\begin{scenariodata}
No additional Fortification counters.\\
\end{scenariodata}

\begin{scenariodata}
\textbf{Confederates:}\\
\end{scenariodata}

\begin{scenariodata}
Resource Factor 214\\
\end{scenariodata}

\begin{scenariodata}
\textbf{Initial Placement:}\\
\end{scenariodata}

\begin{scenariodata}
Hex 0-26 : 141, 1A, (Ramseur)\\
Hex N-26 : 101, 5A, (Kershaw)\\
Hex M-25 : 61, 4A, (Early)\\
\end{scenariodata}

\begin{scenariodata}
Hex 0-25 : 151, 1A, 1C, (Gordon)\\
Hex L-24 :5C, 1H\\
\end{scenariodata}

\begin{scenariodata}
Hex N-27 : 11, (1S, 3W)\\
\end{scenariodata}

\begin{scenariodata}
Hex 0-51 (Staunton) : 1, (1F, 1S, 2W)\\
\end{scenariodata}

\begin{scenariodata}
Hex BB-53 (Charlottesville) : 21, 1A, (1F, 1S)\\
\end{scenariodata}

\begin{scenariodata}
(Partisans: Mosby, Gilmor)\\
\end{scenariodata}

\begin{scenariodata}
\textbf{Withdrawal Hexes:}\\
\end{scenariodata}

\begin{scenariodata}
CC-40 through CC-58\\
\end{scenariodata}

\begin{scenariodata}
\textbf{Supply/Wagon Entry Hexes:}\\
\end{scenariodata}

\begin{scenariodata}
CC-40 through CC-58, 0-51, M-50, E-46\\
\end{scenariodata}

Five additional Fortification counters available.

\begin{scenariodata}
SPECIAL \#1: All Union units in the following hexes must be set up in\\
\end{scenariodata}

Column Formation: N-25, N-24, M-23, N-23, M-22, M-21.

\begin{scenariodata}
(SPECIAL \#2: The entire Shenandoah Valley has been devastated. The\\
\end{scenariodata}

northern limits of this devastation can be marked by placing Destruction counters in hexes Q-21 and L-19 -- this devastation continues to the

\begin{scenariodata}
south edge of the Mapboard.)\\
\end{scenariodata}

\begin{scenariodata}
SPECIAL \#3: The Confederates get a 3-Victory Point bonus at the end\\
\end{scenariodata}

of the game for each of the following hexes that they occupy: N-25,

\begin{scenariodata}
N-24, M-23, M-22, M-21, L-21, L-20, N-22.\\
\end{scenariodata}

\begin{scenariodata}
\textbf{Victory Conditions:} All other results are a draw.\\
\end{scenariodata}

\begin{scenariodata}
\textbf{ADVANCED GAME:} Union wins if they have 15 or more lead in total\\
\end{scenariodata}

\subsection*{Victory Points.}

\begin{scenariodata}
Confederates win if they have 10 or more lead in\\
\end{scenariodata}

total Victory Points.

\begin{scenariodata}
\textbf{BASIC GAME:} This scenario is not suitable for Basic Game play.\\
\end{scenariodata}

\begin{scenariodata}
\textbf{HISTORICAL RESULTS:} Union victory.\\
\end{scenariodata}

\begin{scenariodata}
\textbf{PLAYTESTER'S COMMENTS:} This scenario is not especially interesting\\
\end{scenariodata}

as a two-player game. It is, however, excellent as a solitaire exercise. The

\begin{scenariodata}
Confederates should try to cause as many casualties as possible with their\\
\end{scenariodata}

initial onslaught, and take ground -- especially the Supply Counter in

\begin{scenariodata}
Hex N-24 (otherwise, the Confederates will be unsupplied during the\\
Union Phase\}. The Confederates must take chances and rely on the\\
\end{scenariodata}

generals to rally the units if disorders occur. The Union Phase, its success or failure, will decide the game, which often goes down to the wire.

The Confederates at first swept irresistably forward, the Rebel Yell shrieking out of the morning mists, and the VII! and XIX Corps being routed out of their camps. The attack finally ran out of steam as it hit the tough veterans of the VI Corps. Sheridan made his famous ride from Winchester, the little man on the big horse waving his ridiculous little hat and hollering, "Come on back, boys, give 'em hell!'' The Army of the Shenandoah rallied, reformed its lines, and, late in the afternoon, drove forward in a final counter-attack. The Confederates were routed and driven from the field. Ramseur being among the stain. It was over, and the Shenandoah would see no more major engagements.

The 1864 Shenadoah campaign, which had begun so brightly and brilliantly for the Confederacy, ended in total disaster with the lovely Valley a blackened ruin, and the proud Confederate veterans defeated and humiliated. Jubal Early had to accept the blame for the unhappy series of events, and ended the war in disgrace. Phillip Sheridan became acknowledged as one of the most important architects of the final Union victory, and went on to further successes.

